\chapter{Methodology}
\label{ch:method}

\section{Overview}
\label{sec:method-overview}

\caster\ is an iterative repair loop with one extra component: a small record
kept for each test case, updated every time that test case is run. Everything
else follows from what is in that record.

The record holds three numbers per test case: how many times the test has been
run in this session, how many of those runs rejected a candidate, and how much
CPU time those runs consumed. From those three numbers the loop derives two
estimates --- how likely the test is to reject the next candidate, and how long
it takes --- and spends them in two places. The \textbf{scheduler} uses both to
decide which test to run first. The \textbf{hint selector} uses the first,
together with how readable the test's input is, to decide which facts go into
the next prompt.

Figure~\ref{fig:pipeline} shows one round.

\begin{figure}[htbp]
\centering
\resizebox{0.70\textwidth}{!}{%
\begin{tikzpicture}[
  stage/.style={draw=bkblue,fill=bkblue!6,rounded corners=2pt,align=center,
                inner sep=6pt,font=\small,minimum height=18mm,text width=32mm},
  store/.style={draw=pendred,fill=pendred!6,rounded corners=2pt,align=center,
                inner sep=6pt,font=\small,minimum height=14mm},
  lbl/.style={font=\scriptsize,align=center,inner sep=2pt},
  arr/.style={-{Stealth[length=2.6mm]},draw=black!70,thick},
  darr/.style={-{Stealth[length=2.6mm]},draw=pendred!85,thick,dashed}
]
% --- the four stages, evenly spaced on one row ---
\node (gen)   at ( 0.0,0) [stage] {\textbf{1. Generate}\\[2pt] one model call\\
                                   produces a\\ candidate patch};
\node (sched) at ( 4.4,0) [stage] {\textbf{2. Schedule}\\[2pt] rank the tests,\\
                                   freeze the order};
\node (val)   at ( 8.8,0) [stage] {\textbf{3. Validate}\\[2pt] run tests in\\
                                   that order, stop\\ at the first rejection};
\node (hint)  at (13.2,0) [stage] {\textbf{4. Select hints}\\[2pt] pick $k$ facts\\
                                   for the next\\ prompt};

\draw[arr] (gen)   -- (sched);
\draw[arr] (sched) -- (val);
\draw[arr] (val)   -- (hint);

% --- the store, centred under stages 2-4 ---
\node (store) at (8.8,-3.4) [store,text width=108mm]
  {\textbf{Repair history store} (Section~\ref{sec:store})\\[2pt]
   per test case $x$:\; runs $n(x)$,\; rejections $k(x)$,\; total CPU $S(x)$};

% --- three symmetric vertical links, one per stage above it ---
\draw[darr] (store.north -| sched) -- node[lbl,right,pos=0.55]{read}  (sched.south);
\draw[darr] (val.south)            -- node[lbl,right,pos=0.45]{write} (store.north -| val);
\draw[darr] (store.north -| hint)  -- node[lbl,right,pos=0.55]{read}  (hint.south);

% --- the round-to-round feedback, centred above the row ---
\coordinate (upR) at ($(hint.north)+(0,9mm)$);
\coordinate (upL) at ($(gen.north)+(0,9mm)$);
\draw[arr] (hint.north) -- (upR)
      -- node[lbl,above,pos=0.5]
         {\textbf{round $t{+}1$:} the selected hints enter the prompt}
      (upL) -- (gen.north);
\end{tikzpicture}}
\caption{One round of the \caster\ loop. The repair history store is written by
validation and read by both the scheduler and the hint selector, which is why
the two mechanisms are one system rather than two independent features.}
\label{fig:pipeline}
\end{figure}

\subsection{Requirements the Design Must Satisfy}

Five rules are binding on the implementation, because breaking any of them
produces a plausible-looking but wrong number rather than a visible error. The
store holds $(n, k, S)$ per test case, and the order is derived from it
\textbf{once per round} and then \textbf{frozen} --- re-ranking mid-round would
make the order depend on results from the same scan, which is not reproducible
from the store alone. The scan \textbf{stops at the first rejection}, so only
tests actually run update $n$ and $S$, and only the rejecting test updates $k$.
A candidate is accepted only after surviving \textbf{every} usable test, so
acceptance is independent of the order (Definition~\ref{def:soundness}). And a
hint is rendered as a \textbf{positive requirement} over a concrete
input/output pair, never as a named prohibition; the prompt never names a test
whose content it does not show.

Two measurement rules accompany them. Cost is read as a
\code{getrusage}\allowbreak\code{(RUSAGE\_CHILDREN)} delta --- CPU time, not
wall clock --- because wall clock on a shared machine is contaminated by
sibling processes and would make the learned cost estimate a function of
machine load rather than of the test. And because \code{getrusage} accounting
is per process, parallelism must be process-per-worker rather than threaded,
or one worker's CPU is attributed to another.

\subsection{What the Mechanism Costs}
\label{sec:method-cost}

A method that reduces one cost is only worth having if it does not quietly
introduce another. It is therefore worth saying precisely what \caster\ adds,
because the answer is close to nothing: \textbf{three numbers per test case,
updated in constant time, with nothing trained and nothing executed that the
loop was not already going to execute.}

Per task, the store is $O(|P|)$ memory --- three numbers per test, a few
hundred bytes for a pool of forty. An update is two increments and an
addition; a round is one sort of at most $|P|$ items, built lazily and
abandoned at the first rejection. There are \textbf{no} extra program
executions (the scheduler reorders runs the loop was already making), \textbf{no}
extra model calls (hint selection only reads the store), \textbf{no} growth of
the prompt with session length (at most $k$ hints in any round, not $O(B)$
transcript lines), \textbf{no} training and \textbf{no} offline profiling pass
--- test costs are measured from the runs the loop pays for anyway.

Two consequences follow. \textbf{The information is free}: the scheduler's
inputs --- which test rejected which candidate, and how long that took --- are
by-products of work the loop performs regardless, which is the property that
made \frtp~\citep{qi2013frtp} cheap enough to adopt. And \textbf{the component
is small and portable}: the store is statistics over test-case identifiers, with
no dependency beyond the standard library, no parsing of the program and no
knowledge of the language, which is why it unit-tests without a benchmark, a
sandbox or a model, and why any loop that can run a test and time it can keep
this record.

\paragraph{One caveat, so the claim is not overstated.} The \emph{experiment} in
Chapter~\ref{ch:experiments} does run an extra full-scan \code{audit} pass over
every patch and every test. That pass exists so the kill matrix and the cost
table are established independently of which configuration ran first, which is
a requirement of \emph{comparing} schedules fairly. It is not part of the
method, and a deployment running \caster\ does not perform it.

\section{The Repair History Store}
\label{sec:store}

\subsection{What Is Stored}

For every test case $x$ in the pool, the store holds three counters:

\begin{center}
\small
\begin{tabular}{@{}l l@{}}
\toprule
$n(x)$ & how many times $x$ has been run in this session\\
$k(x)$ & how many of those runs rejected the candidate\\
$S(x)$ & the total CPU seconds those runs consumed\\
\bottomrule
\end{tabular}
\end{center}

Nothing else is stored, and in particular nothing about the program under test.
The store is therefore pure statistics over test-case names: it can be built,
inspected and unit-tested without a benchmark, a sandbox or a model, which is
what makes the scheduler testable in isolation.

\subsection{The Two Derived Estimates}

\paragraph{The rejection rate.} Treat ``does $x$ reject the next candidate?'' as
a coin flip with unknown bias, and estimate that bias from $k(x)$ successes in
$n(x)$ trials. Using the standard Beta conjugate prior, the posterior mean is
\begin{equation}
K(x) \;=\; \frac{k(x) + \alpha}{n(x) + \alpha + \beta}.
\label{eq:K}
\end{equation}
The prior is doing real work here, because at the start of a session every
test has $n = 0$ and the prior is the \emph{only} thing ranking them.

The default is $\mathrm{Beta}(\tfrac12, 2)$, whose mean is $0.2$ --- above
the measured base rate on this corpus and comfortably below one half. The
obvious alternative, Jeffreys' $\mathrm{Beta}(\tfrac12,\tfrac12)$
\citep{jeffreys1946invariant}, places a never-run test at $K = 0.5$, which
outranks a test with a real record of $2$ rejections in $10$ runs ($0.227$). A
rule that always prefers the unmeasured test to the proven one ends up running
the whole pool, so under Jeffreys the scheduler stops scheduling. Trying new
tests is the job of the exploration term below, not a side effect of the
prior. Both priors ship, and the choice is part of the configuration key.

\paragraph{The cost estimate.} The running mean cost of a test is
$\hat{m}(x) = S(x)/n(x)$. A test that has never been run has no mean, and the
cold-start value chosen for it decides how aggressively the loop tries new
tests. Entering it at zero would make an unrun test infinitely attractive;
entering it at the pool average would make the first round's order essentially
arbitrary. \caster\ enters an unrun test at the \emph{maximum} of the observed
per-test means, which is pessimistic: an unmeasured test must earn its place
through its rejection rate rather than through an optimistic guess about its
cost. Two alternative cold starts, the median and the $75$th percentile, are
kept as sensitivity configurations (Section~\ref{sec:exp-config}).

\subsection{Update Rules}

After each scan, the store is updated as follows. Every test that was actually
run has its $n$ incremented and its measured CPU added to $S$. The one test that
rejected the candidate, if any, has its $k$ incremented. Tests in the unrun
tail of the order are untouched, because they cost nothing and revealed nothing.

One accounting rule is worth stating, because violating it silently corrupts the
cost estimate. A test that \emph{times out} is excluded from the cost model
rather than entered at the timeout value. A timeout is weaker evidence than a
wrong answer --- a correct but slow program times out too --- and entering
timeouts at the clock value would let a handful of them dominate $\hat{m}$ for
the whole pool.

\section{Cost-Aware Validation Scheduling}
\label{sec:schedule}

\subsection{The Ranking Rule}

The scheduler ranks tests by
\begin{equation}
\mathrm{score}(x) \;=\; \frac{K(x)}{\hat{m}(x)}
\label{eq:score}
\end{equation}
and runs them in decreasing score until one of them rejects the candidate.

The intuition is a \emph{rate of return}. $K(x)$ is roughly the probability that
spending $\hat{m}(x)$ seconds ends the round; dividing one by the other gives
expected rejections per CPU second, and greedily taking the best return first is
the natural heuristic for the cost in Equation~\eqref{eq:cost}. Informally: run
the test that is most likely to say ``no'' per second spent, not the test that
is most likely to say ``no'' overall. This is full cost-cognizance in the sense
of \citet{elbaum2001cost}, with the difference that $\hat{m}$ here is measured
online from runs the loop was making anyway rather than supplied in advance.

A natural generalisation is to raise the denominator to a power $c$, which at
$c = 0$ recovers ranking on the rejection rate alone and at $c = 1$ recovers
Equation~\eqref{eq:score}; values in between discount a cost estimate that is
noisy early in a session. The implementation and every number in this report
use $c = 1$. Tuning $c$ on the pilot and reporting the result on the same data
would be a selected maximum, so the sweep belongs to the corpus run
(Section~\ref{sec:future}), where the value can be declared in advance.

\paragraph{Exploration.} With probability $\varepsilon = 0.15$, a position in
the order is filled uniformly at random from the remaining tests instead of
greedily. It is in the report rather than left as an implementation detail
because it is load-bearing. Without it, a task whose first test rejects every
candidate reaches the end of the session having measured exactly one test: the
cost estimate never fills in, so the scheduler has nothing to rank by and the
hint selector has one fact to choose from. Because $\varepsilon$ changes measured numbers it is part of the
configuration key. That has a cost worth recording: the random draws are
seeded from that key, so two configurations differing in one flag also explore
differently, which makes a single-flag ablation noisier than it needs to be.

\subsection{Correcting a Marginal Rate}
\label{sec:overlap}

Equation~\eqref{eq:score} ranks each test on its own record. That is a
\emph{marginal} rate, and under stop-at-first-rejection the quantity that
decides cost is a \emph{conditional} one.

Consider two tests that reject exactly the same candidates. Each has the same
record, so each gets the same $K$, and both sit near the top of the order.
Suppose the first of them runs and the candidate \emph{passes} it. The second
is now nearly worthless: everything it is known to reject has already been
shown not to apply here. The marginal $K$ cannot see this, because it never
looks at \emph{which} candidates a test rejected, only how many.

This is the distinction that separates the \emph{total} and \emph{additional}
families in test prioritisation~\citep{rothermel2001prioritizing}, with a
different notion of coverage: not lines of code, but the set of candidates a
test is known to reject. Write $R(x)$ for the set of candidates test $x$ has
been observed to reject so far in the session, and let $P$ be the tests that
have already run \emph{and passed} on the current candidate. Define
\begin{equation}
\mathrm{overlap}(x \mid a) \;=\; \frac{\lvert R(x) \cap R(a)\rvert}{\lvert R(x)\rvert},
\label{eq:overlap}
\end{equation}
and defer $x$ --- move it to the back of this round's order --- when
$\mathrm{overlap}(x \mid a) \ge \tau$ for some $a \in P$, with $\tau = 0.75$
by default and at most one deferral per test per round.

Two properties matter. It \textbf{defers and never drops}: a deferred test
still runs, it runs last, so an accepting round still executes every test and
Definition~\ref{def:soundness} is untouched. And \textbf{the order costs
nothing to abandon}: it is built one position at a time from the store as it
stood when the round began, and construction stops at the first rejection, so
the unprobed tail is never materialised. Scheduling cannot change \emph{which}
candidates are accepted, so the scheduling claim is a pure cost claim and a
difference in repair rate between two schedules is a soundness bug rather than
a result.

Equation~\eqref{eq:overlap} is also what the hint selector uses to drop
near-duplicate hints (Section~\ref{sec:hints}).

\section{Selected Hints}
\label{sec:hints}

\subsection{The Prompt Is Assembled From Slots}

The prompt is never built by string concatenation at a call site. It is
assembled from eight numbered slots, so that a configuration changes which
slots are filled rather than how a sentence is written, and two configurations
can be compared byte for byte.

\begin{table}[htbp]
\centering
\caption{The prompt slots. S1--S4 and S7 are identical in every configuration,
which is what makes round~$1$ byte-identical across the grid; the hint
mechanism is S5 and S6.}
\label{tab:slots}
\small
\begin{tabularx}{\textwidth}{@{}l X l@{}}
\toprule
Slot & Contents & When\\
\midrule
S1 & Role and input/output contract & always\\
S2--S3 & Program label and current source & always\\
S4 & Worked examples, from the task's own public samples & when it has them\\
\textbf{S5} & \textbf{The contract: $k$ selected requirements} & when the store holds $\ge 1$ proven test\\
\textbf{S6} & \textbf{Progress line: a count and a pass/fail pattern} & with S5\\
S7 & Output format & always, and always last\\
S8 & Sampling instruction & one sensitivity arm only\\
\bottomrule
\end{tabularx}
\end{table}

\subsection{What a Hint Is}

S5 opens with one header line and then lists the selected tests.
Figure~\ref{fig:contract} is the block as the model receives it under the
default disclosure.

\begin{figure}[htbp]
\centering
\begin{minipage}{0.93\textwidth}
\footnotesize
\hrule\vspace{5pt}
\begin{verbatim}
Your patch must produce exactly this output on each of these inputs:

  Input #1 - an estimated 43% of attempted fixes get this input wrong:
    stdin: (verbatim, 2 lines)
<<<
5 3
1 4 2 8 5
>>>
    required: (verbatim, 1 line)
<<<
14
>>>
    your last attempt printed `11`

On the input above, your last attempt scored: fail.
\end{verbatim}
\vspace{-8pt}\hrule
\end{minipage}
\caption[The contract block, as the model receives it]{Slots S5 and S6 for one
selected test, at the default disclosure. Everything here is a statement about
what the patch must do or about what was observed; nothing names a behaviour
to avoid, and nothing names the test itself.}
\label{fig:contract}
\end{figure}

Four properties of this form matter, and each is a rule in the implementation.

\paragraph{No instruction names a behaviour to avoid.} The hint says what the
patch must do, never what it must not. This is sharper than ``prefer positive
phrasing'': the failure mode is specifically an \emph{instruction together
with a named target}. On the frozen log, a failure class the model would
otherwise emit $5.4\,\%$ of the time was emitted $14.7\,\%$ of the time once
the prompt named it --- a $2.7\times$ priming inversion --- while the
instruction cut class repetition only from $82.1\,\%$ to $70.9\,\%$ and grew
the prompt from $462$ to $1{,}263$ tokens. A \emph{description} is not an
instruction, which is why \code{your last attempt printed \`{}11\`{}} is
allowed: it states an observation and names no class. Format constraints in S7
(``nothing else'', ``no prose'') stay, because they are identical in every
configuration and changing them would break round-$1$ identity.

\paragraph{Only proven tests are shown.} A test whose estimate comes entirely
from the prior is a guess, and by the mechanism above a wrong guess does not
merely waste tokens, it primes. A test must have rejected at least one
candidate this session before it can appear.

\paragraph{The number shown is the posterior rate, and it is worded as an
estimate.} The figure is $K(x)$ of Equation~\eqref{eq:K} --- the same
Beta-posterior rate the scheduler ranks by, rounded to the nearest percent. It
is \emph{not} the ranking score and \emph{not} the raw frequency $k(x)/n(x)$.
Under the default prior $K(x) = (k(x)+0.5)/(n(x)+2.5)$, so $43\,\%$ is
displayed both by a test with $k{=}4, n{=}8$ and by one with $k{=}1, n{=}1$,
whose tallies are $50\,\%$ and $100\,\%$. The sentence therefore says
\emph{an estimated $43\,\%$ of attempted fixes get this input wrong} rather
than counting past attempts, because a tally would be false in both cases.
Whether the number is shown at all is a separate configuration
(\code{show\_kill\_rate}): it is the only thing in the prompt that says which
requirement is load-bearing, and it has no analogue in the fact-selection
literature, so it is priced rather than assumed.

\paragraph{A test's identifier is never shown.} A file name carries nothing
the model can act on, and on this benchmark it carries something it should
\emph{not} have: contest naming marks the public sample cases, which are the
small, highly discriminating ones. Because no prompt names a test, the rule
that a prompt never names a test whose content it omits holds by
construction.

\subsection{How Hints Are Selected}

At every round the selector ranks proven tests by
\begin{equation}
\mathrm{hintscore}(x) \;=\; K(x) \times \mathrm{legibility}(x)
\label{eq:hintscore}
\end{equation}
and walks the ranking greedily, taking at most $k$ --- with $k = 3$ by
default.

\paragraph{Cost is deliberately absent.} The scheduler and the selector spend
different currencies: a scan pays CPU seconds, a prompt pays tokens. A test
taking five seconds is expensive for the scheduler and free for the prompt, so
carrying $\hat{m}$ into Equation~\eqref{eq:hintscore} would discard the most
informative facts for nothing.

\paragraph{Legibility replaces it.} $\mathrm{legibility}(x)$ is $1$ when the
test's input is at most $400$ characters and $0.45$ otherwise. The threshold
exists because an input of several megabytes cannot be shown, and showing a
truncated prefix of it while asserting a required output produces a
\emph{false fact} --- the model is told to produce a specific output on an
input it has not seen. The incident that produced the rule: the highest
kill-rate test for one task is a $2.3$\,MB stress input, which is an excellent
\emph{schedule} entry and worthless as a \emph{fact}. An illegible input is
therefore rendered as a \emph{shape} --- its line count, its character count,
and its first line verbatim, because on this benchmark the first line almost
always carries $N$ and so says whether the failure is a complexity problem or
an edge case. An expected output too large to show is \emph{withheld} rather
than shaped, since a required value the model cannot read is noise with an
assertion attached. The penalty is a factor rather than a veto, so a dominant
illegible test can still outrank a legible one when nothing else is close.

\paragraph{Near-duplicate hints are dropped.} Two hints that reject the same
candidates are one fact spending two slots. The selector skips a candidate
hint whose rejection set is already covered by one it has chosen, using the
containment of Equation~\eqref{eq:overlap} and the same threshold $\tau$ ---
but \emph{symmetrically}, unlike in the scheduler:
\begin{equation}
\max\bigl(\mathrm{overlap}(x \mid m),\; \mathrm{overlap}(m \mid x)\bigr) \;\ge\; \tau
\quad\text{for some already-chosen } m \;\Longrightarrow\; \text{skip } x .
\label{eq:hintdedup}
\end{equation}
In the scheduler one test has already run and the question is whether the next
still adds anything; here neither has run and the question is whether the
\emph{pair} is worth two slots, so a narrow hint inside a broad one is dropped
whichever way round the two are ranked.

\paragraph{$k$ is a ceiling, not a count.} Between the deduplication above and
a thin pool of proven rejecters --- $55\,\%$ of pilot sessions ended with
fewer than two --- the selector returns fewer than $k$ far more often than it
returns $k$. Every record therefore carries the number of hints that actually
reached the prompt alongside the number requested, and the sweep of
Chapter~\ref{ch:experiments} is read against the former. A sweep read against
the requested $k$ would report a comparison it did not run.

\subsection{The Disclosure Dial}

How much of a test case the prompt reveals is a parameter, not a fixed choice,
because it trades two risks against each other. Showing less gives the model too
little to work with. Showing more invites it to satisfy \emph{the shown tests}
specifically --- writing a special case for the displayed input rather than
repairing the fault --- which produces a patch that passes the hinted tests and
fails everything else.

\caster\ therefore exposes three levels, measured as a dose-response curve in
Chapter~\ref{ch:results}:

\begin{table}[htbp]
\centering
\caption{Disclosure levels, from least to most revealed.}
\label{tab:disclosure}
\small
\begin{tabularx}{\textwidth}{@{}l X@{}}
\toprule
Level & What the prompt shows\\
\midrule
\code{sample-only} & \textbf{Nothing from the store.} The contract block is
withheld in full, so the only worked examples in the prompt are the task's own
public samples. This is the floor of the ladder, not a partial disclosure.\\
\code{no-required-line} & The selected \emph{inputs}, and a line saying the
patch is checked on them --- but their expected outputs are withheld.\\
\code{full} (default) & The selected inputs, their expected outputs, and an
explicit statement that the patch must produce them.\\
\bottomrule
\end{tabularx}
\end{table}

The number of hints $k$ and the disclosure level are the same dial seen from two
sides, and Chapter~\ref{ch:results} reports them as one curve rather than two
findings.

\subsection{Why This Is Not a Transcript}

Three differences follow from the construction. \emph{The prompt does not
grow}: at most $k$ hints appear in any round and the bound does not depend on
how long the session has run, where a transcript grows linearly in rounds and
therefore walks past the point where more context starts to
hurt~\citep{parasaram2025factselection}. \emph{The content changes}: which
tests have the highest rejection rate moves as the session proceeds, so round
$10$ can carry entirely different hints from round $3$, where a transcript is
append-only. \emph{The facts are earned, not asserted}: each hint is backed by a measured
record produced by the scheduler that had to run the test anyway, and the hint
mechanism adds no executions of its own. Even the one line that reports the
previous attempt carries the \emph{last} outcome only --- an accumulating
record grows the prompt roughly $3\times$ for a benefit measured at zero.
